\begin{figure*}
    \centering
    \includegraphics[width=2.57cm]{./images/brasao_da_republica.jpeg}
    
    \textsf{Ministério da Educação}\\
    \textsf{Universidade Federal do Agreste de Pernambuco}\\
    \textsf{Bacharelado em Ciência da Computação}
\end{figure*}

\vspace{0.5cm}

\begin{scriptsize}
\noindent \fbox{
\begin{minipage}[t]{\textwidth}
\textsf{\scriptsize COMPONENTE CURRICULAR:} \\
\textbf{DESENVOLVIMENTO E EXECUÇÃO DE PROJETOS DE SOFTWARE} \\
\textsf{\scriptsize CÓDIGO:} \textbf{UAG00305}
\end{minipage}}

\noindent \fbox{
\begin{minipage}[t]{0.293\textwidth}
\textsf{\scriptsize PERÍODO A SER OFERTADO:} \\
\textbf{0}
\end{minipage}
%
\begin{minipage}[t]{0.7\textwidth}
\textsf{\scriptsize NÚCLEO DE FORMAÇÃO:} \\
\textbf{COMPONENTES OPTATIVOS ÁREA TEMÁTICA METODOLOGIA E TÉCNICAS DA
COMPUTAÇÃO}
\end{minipage}}

\noindent \fbox{
\begin{minipage}[t]{0.493\textwidth}
\textsf{\scriptsize TIPO:} \textbf{OPTATIVO}
\end{minipage}
%
\begin{minipage}[t]{0.5\textwidth}
\textsf{\scriptsize CRÉDITOS:} \textbf{4}
\end{minipage}}

\noindent \fbox{
\begin{minipage}[t]{0.3\textwidth}
\textsf{\scriptsize CARGA HORÁRIA TOTAL:} 
\textbf{60}
\end{minipage}
%
\begin{minipage}[t]{0.19\textwidth}
\textsf{\scriptsize TEÓRICA:} 
\textbf{30}
\end{minipage}
%
\begin{minipage}[t]{0.19\textwidth}
\textsf{\scriptsize PRÁTICA:} 
\textbf{30}
\end{minipage}
%
\begin{minipage}[t]{0.30\textwidth}
\textsf{\scriptsize EAD-SEMIPRESENCIAL:} 
\textbf{0}
\end{minipage}}

\noindent \fbox{
\begin{minipage}[t]{\textwidth}
\textsf{\scriptsize PRÉ-REQUISITOS:}
CCMP3057 INTRODUÇÃO À PROGRAMAÇÃO
\end{minipage}}

\noindent \fbox{
\begin{minipage}[t]{\textwidth}
\textsf{\scriptsize CORREQUISITOS:} 
Não há.
\end{minipage}}

\noindent \fbox{
\begin{minipage}[t]{\textwidth}
\textsf{\scriptsize REQUISITO DE CARGA HORÁRIA:} 
Não há.
\end{minipage}}

\noindent \fbox{
\begin{minipage}[t]{\textwidth}
\textsf{\scriptsize EMENTA:} \\
Práticas em tecnologias e frameworks para desenvolvimento de aplicações
Web. Preparação de ambientes reais para especificação, desenvolvimento e
implantação de sistemas de software. Definição e implantação de
processo/metodologia de desenvolvimento de software em projetos reais.
\end{minipage}}

\noindent \fbox{
\begin{minipage}[t]{\textwidth}
\textsf{\scriptsize BIBLIOGRAFIA BÁSICA:}
\begin{enumerate}
\def\labelenumi{\arabic{enumi}.}
\item
  Horstmann, Cay \& Geary, David. Core Javaserver Faces. Editora Alta
  Books, tradução da 3a edição Americana, 2012.
\item
  Goncalves, Edson. Dominando. Java Server Faces e Facelets Utilizando
  Spring 2.5, Hibernat e JPA. CIENCIA MODERNA; Edição: 1a 2008.
\item
  Wazlawick, Raul. Análise e Design Orientados a Objetos Para Sistemas
  de Informação: Modelagem com UML, OCL e IFML. Ed.: Campos, 2014.
\end{enumerate}
\end{minipage}}

\noindent \fbox{
\begin{minipage}[t]{\textwidth}
\textsf{\scriptsize BIBLIOGRAFIA COMPLEMENTAR:}
\begin{enumerate}
\def\labelenumi{\arabic{enumi}.}
\item
  LUCKOW, D. H \& MELO, A. A. Programação Java Para a Web: Aprenda a
  Desenvolver uma Aplicação Financeira Pessoal com as Ferramentas Mais
  Modernas da Plataforma Java. Novatec; Edição: 2a, 2015.
\item
  CAY S. HORSTMANN E GARY CORNELL. Core Java 2. Volume I - Fundamentos.
  Sun Microsystems Press, Makron Books do Brasil, 2001
\item
  LIMA, Adilson da Silva. UML 2.5: do requisito à solução. São Paulo:
  Ed.: Érica. Edição: 1a, 2014.
\end{enumerate}
\end{minipage}}
\end{scriptsize}

\newpage